% Required packages (add to main.tex preamble if not already present):
%   \usepackage{amsmath}
%   \usepackage{algorithm}
%   \usepackage{algorithmic}
%   \usepackage{booktabs}

\section{Methodology}
\label{sec:method}

\subsection{Problem Formulation}

We consider a continual fine-tuning setting with $K$ tasks arriving sequentially.
Let $\mathcal{T} = \{T_1, T_2, \ldots, T_K\}$ denote the task sequence, where each task $T_k$ is
defined by a training set $\mathcal{D}_k^{\mathrm{train}}$ and an evaluation set
$\mathcal{D}_k^{\mathrm{val}}$.
A shared language model $f_\theta$ parameterized by $\theta$ is adapted to each task in order; upon
receiving $T_k$, the model may only access $\mathcal{D}_k^{\mathrm{train}}$ and, for rehearsal
methods, a bounded memory buffer $\mathcal{M}$ populated from prior tasks.
After training on the full sequence, let $R[i,j]$ denote the evaluation metric of the model on task
$T_j$ immediately after training on task $T_i$ ($i \geq j$), forming the $K \times K$
\emph{performance matrix} $R$.

\paragraph{Evaluation metrics.}
Three continual learning metrics are derived from $R$:
\begin{align}
  \mathrm{AA}  &= \frac{1}{K}   \sum_{j=1}^{K}
                   R[K,\,j],
                   \label{eq:aa} \\
  \mathrm{F}   &= \frac{1}{K-1} \sum_{j=1}^{K-1}
                   \Bigl(\max_{i \leq K-1} R[i,\,j] - R[K,\,j]\Bigr),
                   \label{eq:forgetting} \\
  \mathrm{BWT} &= \frac{1}{K-1} \sum_{j=1}^{K-1}
                   \bigl(R[K,\,j] - R[j,\,j]\bigr).
                   \label{eq:bwt}
\end{align}
Average Accuracy (AA) measures final performance across all tasks (Eq.~\ref{eq:aa}).
Forgetting (F) quantifies the mean drop from peak performance to final performance on each past
task; lower is better (Eq.~\ref{eq:forgetting}).
Backward Transfer (BWT) captures whether learning later tasks helps or hurts earlier ones; positive
BWT indicates beneficial transfer (Eq.~\ref{eq:bwt}).

\subsection{Baselines}

\paragraph{Full Fine-Tuning (Full-FT).}
All model parameters $\theta$ are updated on each task in sequence with no continual learning
mechanism.
Full-FT establishes an upper bound on single-task performance and a lower bound on retention of
prior tasks; it serves as the na\"{i}ve baseline.

\paragraph{Vanilla LoRA.}
A single shared LoRA adapter $(A, B)$ is trained across the entire task sequence, with the backbone
frozen.
No mechanism prevents interference between task-specific gradient updates; this represents the
minimal PEFT baseline and quantifies the forgetting attributable to sequential adapter training
alone.

\paragraph{LoRA + EWC.}
Elastic Weight Consolidation~\cite{kirkpatrick2017overcoming} augments the per-task loss with a
quadratic penalty that resists deviation of trainable parameters from their values at the end of
the previous task:
\begin{equation}
  \mathcal{L}_{\mathrm{EWC}}(\theta)
    = \mathcal{L}_k(\theta)
    + \frac{\lambda}{2} \sum_{i}
        \hat{F}_i \bigl(\theta_i - \theta_i^*\bigr)^2,
  \label{eq:ewc}
\end{equation}
where $\hat{F}_i$ is the $i$-th diagonal element of the empirical Fisher information matrix
computed over a held-out subset of $\mathcal{D}_{k-1}^{\mathrm{train}}$, $\theta^*$ are the
adapter parameters after training on $T_{k-1}$, and $\lambda > 0$ is a task-independent scaling
coefficient.
Only LoRA adapter parameters participate in the EWC penalty; the frozen backbone is excluded.

\paragraph{LoRA + Replay (ER).}
A fixed-size episodic buffer $\mathcal{M}$ stores a random subset of examples from each completed
task.
At each training step on $T_k$, a mini-batch is formed by concatenating examples from
$\mathcal{D}_k^{\mathrm{train}}$ with a uniform sample from $\mathcal{M}$.
No regularization term is added; this method isolates the effect of raw rehearsal.

\paragraph{O-LoRA.}
Orthogonal LoRA~\cite{wang2023olora} learns task-$k$ adapters in the subspace orthogonal to all
prior tasks by minimizing
\begin{equation}
  \mathcal{L}_{\mathrm{O\text{-}LoRA}}(\theta)
    = \mathcal{L}_k(\theta)
    + \mu \sum_{j < k} \left\| A_k^\top A_j \right\|_F^2,
  \label{eq:olora}
\end{equation}
where $A_k, A_j \in \mathbb{R}^{d \times r}$ are the LoRA $A$-matrices for tasks $k$ and $j$
respectively, and $\mu > 0$ controls orthogonality strength.
The prior-task matrices $A_j$ are frozen after their respective tasks complete and stored for use
in subsequent penalty computations.

\subsection{Minimal Mixed Rehearsal (MMR)}
\label{sec:mmr}

We propose \textbf{Minimal Mixed Rehearsal (MMR)}, a conceptually simple strategy that combines
experience replay with QLoRA's parameter efficiency without any additional regularization or
architectural constraints.
The central design choice is a small \emph{replay ratio} $f \in (0,1)$: at each optimization step,
a fraction $f$ of the mini-batch is drawn from the replay buffer $\mathcal{M}$ and the remaining
$(1-f)$ from the current task.
After completing task $T_k$, exactly $B$ examples are sampled uniformly from
$\mathcal{D}_k^{\mathrm{train}}$ and appended to $\mathcal{M}$.

Algorithm~\ref{alg:mmr} provides a complete specification of MMR.

\begin{algorithm}[t]
\caption{Minimal Mixed Rehearsal (MMR)}
\label{alg:mmr}
\begin{algorithmic}[1]
\REQUIRE Model $\theta$, task sequence $\{T_1,\ldots,T_K\}$, replay ratio $f \in (0,1)$,
         buffer capacity $B$ per task, mini-batch size $M$, learning rate $\eta$
\ENSURE  Trained model $\theta$
\STATE Initialize replay buffer $\mathcal{M} \leftarrow \emptyset$
\FOR{$k = 1$ \textbf{to} $K$}
  \STATE $n_{\mathrm{replay}} \leftarrow \lfloor f \cdot M \rfloor$;\quad
         $n_{\mathrm{curr}}   \leftarrow M - n_{\mathrm{replay}}$
  \FOR{each mini-batch $\mathbf{b}_{\mathrm{curr}} \sim \mathcal{D}_k^{\mathrm{train}}$
       of size $n_{\mathrm{curr}}$}
    \IF{$|\mathcal{M}| > 0$}
      \STATE $\mathbf{b}_{\mathrm{replay}} \sim \mathcal{M}$,\quad size $n_{\mathrm{replay}}$
      \STATE $\mathbf{b} \leftarrow \mathbf{b}_{\mathrm{curr}} \,\|\,
             \mathbf{b}_{\mathrm{replay}}$ \COMMENT{Concatenate batches}
    \ELSE
      \STATE $\mathbf{b} \leftarrow \mathbf{b}_{\mathrm{curr}}$
    \ENDIF
    \STATE $\theta \leftarrow \theta - \eta\,\nabla_\theta\,\mathcal{L}(\theta;\,\mathbf{b})$
  \ENDFOR
  \STATE $\mathcal{M} \leftarrow \mathcal{M} \cup
         \mathrm{UniformSample}(\mathcal{D}_k^{\mathrm{train}},\,B)$
         \COMMENT{Buffer update}
\ENDFOR
\RETURN $\theta$
\end{algorithmic}
\end{algorithm}

MMR introduces \emph{no new trainable parameters} and requires no Fisher information computation.
Its memory overhead is $O(K \cdot B)$ stored examples, which is negligible when $B \ll
|\mathcal{D}_k^{\mathrm{train}}|$.
The replay ratio $f$ is the sole additional hyperparameter; we study its sensitivity in
Section~\ref{sec:results}.
Compared with LoRA + Replay (ER), MMR differs in enforcing a strict fraction-based mixing policy
at every step (rather than concatenating full buffer batches), which produces a more controlled
trade-off between current-task loss and retention loss.

\subsection{Implementation under QLoRA}

All methods are implemented within the QLoRA~\cite{dettmers2023qlora} framework.
The backbone is loaded in 4-bit NF4 precision using \texttt{BitsAndBytesConfig} from the
\texttt{bitsandbytes} library, with \texttt{bfloat16} as the compute dtype to leverage Tensor Core
throughput on the NVIDIA T4.
We use the \texttt{paged\_adamw\_8bit} optimizer, which pages optimizer states to CPU when GPU
memory is exhausted, enabling training within the 16\,GB budget.
LoRA adapters are attached to all four attention projection matrices
(\texttt{q\_proj}, \texttt{k\_proj}, \texttt{v\_proj}, \texttt{o\_proj}) with rank $r{=}16$,
scaling $\alpha{=}32$, and dropout $p{=}0.05$.
For O-LoRA, we additionally store the frozen $A$-matrices from each completed task to compute the
orthogonality penalty in Eq.~\ref{eq:olora}.
PEFT model management uses the \texttt{peft} library~\cite{mangrulkar2022peft}.
